\chapter{Dual-Process Accounts of Deductive Reasoning}
\label{sec:p19-rationality-2-5}

Klauer's 2021 chapter distinguishes cognitive accounts that are often
grouped under ``dual process,'' compares their functions and interactions,
and assesses evidence for and against them. A correct answer, a long
response time, or a believable conclusion does not by itself identify
a cognitive process \cite[pp.~173--184]{Klauer2021DualProcess}.
This treatment uses the Evans--Stanovich default-interventionist account
as an explicit reference account and keeps its rivals separate. It is
not a definition making Type~1 the wrong answer and Type~2 the right one.

\section{Process distinctions and interaction architectures}

\begin{hbdefinition}{def:p19-r25-process-account}{What a dual-process account must specify}
A dual-process account postulates qualitatively different kinds of
processing that interact or compete in producing responses. It must
identify the distinguishing features, the information and representations
each process uses, its computed function or proposed response, the
conditions of engagement, and the mechanism resolving their interaction.
Merely naming two response categories does not supply such an account.

In the chapter's initial framework, Type~1 processes encode premises,
direct stimulus-driven attention, activate background knowledge and
pragmatic interpretations, and often propose an intuitive answer. They
can also supply a fallback when deliberation produces no answer.
Type~2 processes implement instructions and a task set, monitor progress
against goals, allocate attention, and support working-memory-dependent
simulation and deliberation. Conflict, absence of an answer, surprise,
or threat can prompt their engagement. These functions are hypotheses
about mechanisms, not a universal ordering of accuracy or usefulness.
\end{hbdefinition}

The source's table lists nine commonly proposed contrasts. They are
retained below as proposed dimensions, not nine coextensive definitions
\cite[table~2.5.1, p.~174]{Klauer2021DualProcess}.
\begin{center}
\begin{tabular}{ll}
\toprule
Proposed Type~1 attribute & Proposed Type~2 attribute\\
\midrule
Fast & Slow\\
Automatic & Controlled\\
Nonconscious & Conscious\\
Unintentional & Intentional\\
Associative & Reflective\\
Parallel & Serial\\
Heuristic & Analytic\\
Experiential & Rational\\
Does not require working memory & Requires working memory\\
\bottomrule
\end{tabular}
\end{center}
The last-but-one row records the source's historical terminology;
it does not establish that experiential processing is irrational.
The later argument explicitly questions treating speed, automaticity,
resource use, and correctness as interchangeable classifications.

\begin{hbdefinition}{def:p19-r25-interactions}{Default intervention, parallel competition, and representation shaping}
In a \emph{default-interventionist} architecture an initial proposal
is acted upon unless another process intervenes, confirms it, or replaces
it. A \emph{parallel-competitive} architecture permits both processes
to generate proposals concurrently and compete for the eventual response.
Neither architecture is identical to a response-time label.

A separate interaction is \emph{representation shaping}: initial attention,
pragmatics, or retrieved knowledge determines the representation on which
subsequent operations act. A demanding reasoning procedure can therefore
operate on a restricted or misleading representation. The source treats
representation shaping and default intervention as compatible interactions,
not mutually exclusive alternatives.
\end{hbdefinition}

\section{Stanovich's tripartite distinctions and failure modes}

Stanovich's framework, as reconstructed by Klauer, calls the heterogeneous
autonomous processes collectively Type~1. They include preattentive
selection, domain-specific modules, implicit learning and conditioning,
and inference or classification rules made automatic by practice or
available without such practice. These are not asserted to form one
homogeneous module \cite[p.~174]{Klauer2021DualProcess}.

\begin{hbdefinition}{def:p19-r25-tripartite}{Autonomous, algorithmic, and reflective functions}
The \emph{algorithmic} level implements strategies, maintains hypothetical
representations, and can inhibit an initial proposal. The \emph{reflective}
level concerns control states and goals that initiate and guide such
operations, including a decision to reconsider or override a proposal.
Thus ability to perform a computation and disposition to initiate it are
different explanatory variables. The source motivates this distinction
both by those functions and by individual-differences research; it does
not turn every observed correlation into a causal separation.
\end{hbdefinition}

The chapter preserves four different failure explanations:
\begin{enumerate}
\item Deliberation is not engaged: the initial answer feels compelling,
      resources are unavailable, or effort is not allocated.
\item Deliberation is engaged but fails to override: the relevant skill
      or rule may not have been learned, a \emph{mindware gap}, rather
      than a simple refusal to think.
\item The available rule or belief system is itself faulty:
      \emph{contaminated mindware} can make extended reasoning biased.
\item Deliberation operates on a focal representation that omits an
      important alternative. In the selection task, initially attending
      only to the symbols mentioned by the rule can lead to a subsequent
      rationalization of the wrong card choice.
\end{enumerate}
These explanations are not identified by a wrong answer alone. They
suggest different tests of engagement, knowledge, representation, and
control; increasing response time need not repair all four.

\section{Tasks, benchmarks, and rationality}

The chapter's figure presents three distinct tasks, not three interchangeable
measures of an internal system \cite[figure~2.5.1, p.~175]{Klauer2021DualProcess}.
\begin{hbexample}{def:p19-r25-task-examples}{Selection, belief conflict, and a reduced base-rate problem}
In the selection task, cards show \(A,B,3,1\); each has a capital
letter on one side and a digit on the other. The rule is that a card
with \(A\) has \(3\) on its opposite side. One must choose cards
to test the rule, not to find more confirming examples.

In the belief task, instructions require assuming the premises true
even when they conflict with reality. The source's training example
stipulates that finishing a drink makes it full, then states that the
drink is finished, and asks whether it is empty. Its trial stipulates
that a crying child is happy, then states that the child is crying and
asks whether the child is happy. The ordinary background expectation
about crying conflicts with the supplied conditional.

In the reduced base-rate task, a person is randomly drawn from a
sample of 999 nannies and five politicians, and described as dishonest.
The response is a category judgment. Counting the larger category
and matching a stereotypical description can both be simple, fast
procedures; their information contents do not fix their processing type.
These are task stipulations, not measured occupational trait frequencies.
\end{hbexample}

\begin{hblemma}{thm:p19-r25-benchmarks}{proved=annex:p19.r25.benchmarks}{The stipulated task benchmarks and the missing likelihood}
Under the selection task's classical rule, the cards requiring inspection
are \(A\) and \(1\), not \(A\) and \(3\). Under the belief task's
stipulated premises, happiness follows by modus ponens irrespective of
ordinary plausibility. These validity results do not specify a participant's
representation or process.

For the reduced category example, if \(D\) is the dishonest description
and its two positive likelihoods are supplied, the posterior odds of
politician to nanny are
\[
 \frac{p(\mathrm{politician}\mid D)}{p(\mathrm{nanny}\mid D)}
   =\frac5{999}\,
      \frac{p(D\mid\mathrm{politician})}{p(D\mid\mathrm{nanny})}.
\]
The description reverses the base-rate ranking only if its likelihood
ratio exceeds \(999/5=199.8\). Without those likelihoods, the base
counts and a stereotype alone do not specify a unique full Bayesian
posterior. A researcher's ``base-rate-consistent'' response label must
not be silently equated with that posterior calculation.
\end{hblemma}

Klauer distinguishes \emph{instrumental rationality}, suitable means to
one's ends, \emph{normative rationality}, adherence to specified norms,
and \emph{epistemic rationality}, coherent and consistent beliefs supported
by evidence. Autonomous stimulus-triggered processes do not directly
read off a newly specified goal, but their trained or environmentally
attuned responses can be correct and useful. A well-practiced modus
ponens response may be fast; a fast-and-frugal procedure may be successful
in a suitable environment. Conversely, deliberate computation can
rationalize a focal belief or apply faulty rules. Instrumental success
and normative correctness therefore do not by themselves divide the
two process types \cite[pp.~174--176]{Klauer2021DualProcess}.

The source separately reports Evans's two-minds conception: the newer
mind, operating on explicit knowledge, pursues epistemic rationality,
whereas old and new minds contribute together to instrumental rationality.
This is that particular account's assignment of functions, not an
empirically established identity between explicit deliberation and true
belief. Klauer uses it to motivate the subsequent justification question.

The chapter also considers \emph{justification} as a possible qualitative
criterion. A correct response arising from an environmental correlation
can fail when the correlation is removed. Requiring awareness of reasons
while generating the response, rather than only a post-hoc explanation,
may distinguish reflective inference from an intuition. Its analogy
with justified true belief is a conceptual argument, not an experiment
showing that every intuitive answer lacks justification. Being able to
give reasons also does not guarantee that the reasons are sound.

\section{Four distinct source accounts}

\subsection{The reference default-interventionist framework}

In the Evans--Stanovich framework, Type~2 processing requires working
memory and supports cognitive decoupling: keeping a supposition separate
from an existing belief while simulating its consequences. Type~1
processing does not require that working-memory support. An intuitive
default may be overridden, or affirmed, depending on motivation,
available capacity, and thinking dispositions. The reflective level
regulates intervention and the algorithmic level implements it. Klauer
emphasizes that this is a metatheory capable of inspiring task models,
not an already complete model of every trial's latency or response
\cite[p.~176]{Klauer2021DualProcess}.

To obtain the familiar belief-bias prediction one must add a task-specific
assumption: retrieving the conclusion's plausibility is easier and less
resource demanding than checking its logical status. Under that assumption,
time pressure or executive load should selectively impair conflict
resolution, whereas a congruent belief can already supply the correct
answer. If plausibility assessment instead requires integration with a
minor premise, this additional assumption is no longer automatic. That
is a discriminating task condition, not permission to rename whichever
process was faster after seeing the result.

\subsection{Oaksford and Chater: one function, two modes of access}

The single-function account of conditional reasoning uses a long-term-memory
constraint-satisfaction network. Encountering a conditional activates
antecedent, consequent, and associated representations. A minor premise
clamps the corresponding node; activation settles and the consequent's
activation is read as a probability assessment. Type~2 queries can be
strategic, involve hypothetical or counterfactual premises, store and
integrate the outputs in working memory, undo effects of prior queries,
and shield a hypothetical representation from irrelevant intrusion.
The representations can resemble mental models, and screened processing
can approximate logical thought as a limiting case, but both modes serve
the same function: assessing probabilities
\cite[pp.~176--177]{Klauer2021DualProcess}.

Klauer presents nonconscious automatic spreading activation and intentional
working-memory querying as plausible features of this account, explicitly
noting that some are his interpretation of its architecture rather than
claims its authors state in precisely those terms. Initial activation
shapes what is considered; strategic querying can modify or inhibit the
default assessment. It is therefore not a rival that identifies intuition
with probability and deliberation with a wholly different formal target.

\subsection{Dual strategy: likelihood versus a counterexample}

Verschueren and colleagues' dual-strategy account, developed by Markovits
and colleagues, separates a statistical assessment from counterexample
search. The statistical strategy accepts a sufficiently probable conclusion
and can be fast and inexpensive. Given time and capacity, a counterexample
strategy can replace it: reject if at least one defeating case is found.
The defeating case can come from the supplied information or strategic
retrieval of background knowledge. Thus the computed functions differ:
frequency or likelihood matters to the first, whereas existence of even
one counterexample is decisive to the second
\cite[p.~177]{Klauer2021DualProcess}.

A rare counterexample can therefore discriminate the strategies in a
way that merely measuring speed cannot. This test needs the assumed
interpretation of the conditional and the availability of the counterexample
specified; inability to retrieve a case is not evidence that none exists.
Unlike the single-function account, this is explicitly a dual-function
proposal. Neither a rare counterexample's relevance to deduction nor a
likelihood-based response's usefulness is settled by giving it a Type label.

\subsection{Dual source: form and background knowledge}

Klauer, Beller, and H\"utter's dual-source model of probabilistic conditional
reasoning distinguishes proposals derived from argument form and from
background knowledge. It is not committed to aligning these sources
with every automatic/controlled, fast/slow, or working-memory contrast.
Its purpose is to make their separate contributions experimentally
accessible, so those further features can be studied rather than assumed
\cite[p.~177]{Klauer2021DualProcess}.

One must accordingly vary or measure the content-derived proposal
separately from the argument form and test their interaction. An inferred
source contribution does not, on its own, identify an autonomous or
deliberate implementation. The selected chapter supplies this architecture
and methodological purpose but no full set of the external model's
equations or fitting procedures. A quantitative test additionally needs
those equations, their parameter constraints, and identifiability checks.

\section{Discriminating predictions and their evidence}

\begin{center}
\begin{tabular}{p{0.25\textwidth}p{0.33\textwidth}p{0.33\textwidth}}
\toprule
Diagnostic contrast & Specified prediction & Rival or limiting explanation\\
\midrule
Load on belief--logic conflict & Under the reference task mapping,
conflict accuracy deteriorates more than congruent accuracy. & A task-specific
resource difference need not uniquely identify two global systems.\\
Plausibility versus validity instructions & If plausibility is an easy
default, it should interfere with validity more than the reverse. & If both
require comparable premise integration, interference can be symmetric.\\
Rare defeating case & A counterexample strategy rejects when the case
is found, even if the conclusion is highly probable. & A single probability
function need not treat existence alone as decisive.\\
Response time by response category & A specified latency model may
predict an ordering. & Choice strength, difficulty, switching, reading, and
response entry can change that ordering without changing architecture.\\
Initial/final response & An intervention account permits correction,
confirmation, or failed reconsideration. & Stable strategic choice or two
initially available easy rules can produce unchanged answers.\\
\bottomrule
\end{tabular}
\end{center}

\subsection{A fully specified primary load study}

\begin{hbdefinition}{def:p19-r25-deneys-record}{De Neys's 2006 executive-load experiment}
The published primary studied 308 first-year psychology students at
the University of Leuven, Belgium, participating for course credit.
They completed a group-adapted operation-span measure, then eight
syllogistic validity problems: four belief--logic conflicts and four
congruent problems. Instructions required treating the premises as true.
Participants were randomly assigned to no, low, or high load. In the
load conditions they memorized a \(3\times3\) dot pattern shown for
850~ms before each syllogism and reproduced it after answering. Low
load used a simple three-dot horizontal pattern; high load used complex
four-dot patterns. The responses were typed accept/reject judgments,
scored by the stated logical benchmark
\cite[pp.~429--430]{DeNeys2006DualProcessing}.

Span groups were the upper, middle, and lower thirds of measured capacity,
not randomly assigned cognitive types. The analysis crossed three load
conditions and three span groups with within-person conflict status.
The adult university sample does not represent children, clinical
populations, or all reasoners. Eight trials per person are not 2,464
independently recruited participants.
\end{hbdefinition}

The primary reports a load effect on conflict accuracy,
\(F(2,299)=7.64,\eta_p^2=.05\), and a linear load trend,
\(F(1,299)=15.28,\eta_p^2=.05\). It reports no corresponding
no-conflict load effect, \(F(2,299)<1\), and no significant load-by-span
interaction, \(F(4,299)<1\). Its inferential convention includes
\(p_{\mathrm{rep}}\); that quantity must not be printed as an ordinary
\(p\)-value. Confidence intervals for these effects are not supplied.
Nonsignificance is not proof of a zero effect or exact group equivalence
\cite[pp.~430--431]{DeNeys2006DualProcessing}.

The groupwise high-load conflict decreases reported for high, middle,
and low spans were 13\%, 16\%, and 22\%, respectively; their difference
was not significant. The dot recall check found about 89\% of dots
correct in the complex condition and 94\% in the simple condition.
These measures concern the manipulation and task, not direct counts
of an unobserved process. Complexity and dot number changed together.

The decisive rival prediction was not simply ``two versus one process.''
A strong individual-difference hypothesis proposed that low-span
reasoners lacked analytic access and should therefore show no further
conflict decrement under load; a weaker hypothesis proposed common
access with different success rates and predicted decrements in both
groups. The observed pattern challenges the strong hypothesis in this
sample. The paper's note~1 explicitly says its comparison does not
settle whether both mechanisms can be unified under a single system.
Its discussion also offers an argument that shared attempts to follow
logic support that norm; observed adherence, even if shared, does not
by itself establish a normative proposition. The task instructions
supply a local validity standard, not a proof that every real-world
reasoning goal is deductive \cite[pp.~429,432--433]{DeNeys2006DualProcessing}.

\subsection{Source-reviewed speed, instruction, and switching contrasts}

Klauer reports that Evans and Curtis-Holmes's time-pressure study
reduced logical discrimination and increased belief bias, whereas
De Neys's load study reduced discrimination without increasing the
belief effect. These are different manipulations and outcomes, not
one generic effect of thinking less. He then contrasts Handley,
Newstead, and Trippas's five experiments: in some conditions belief
judgments were slower than validity judgments, and logical status
interfered with belief judgments \cite[pp.~178--179]{Klauer2021DualProcess}.

The actual stimulus comparison matters. An already believable or
unbelievable conclusion about astronauts and health can prompt a simple
memory assessment. In Handley's first three experiments, ``the child
is happy'' acquired a plausibility assessment only after integration
with the minor premise that the child was crying and the relevant
background conditional. Judging validity integrated the same minor
premise with the experiment's supplied conditional. The operations
could therefore be closer in demand than the usual task mapping assumes.
For Handley's later experiments, and for later alternating-instruction
studies, Klauer retains switching aftereffects as an alternative source
of interference; switching can increase crosstalk between tasks.

The source makes a similar argument for Chun and Kruglanski's base-rate
studies: ease of processing, rather than whether information is a base
rate or stereotype, can influence its weight. The reduced base-rate
format makes both information sources short and can support two easy
rules. Trippas and colleagues' more difficult logical problems reduced
logic's interference with belief judgments, but the chapter also notes
their switching qualification. The selected chapter does not report
all these studies' sample sizes, recruitment, effect intervals, or
switching distributions. They remain explicitly source-reviewed findings,
not independent primary-data estimates supplied by this entry.

\subsection{Response time and the competing critique}

Krajbich and colleagues' critique, as discussed by Klauer, observes that
a preferred response can be faster as its preference becomes stronger.
A sequential-sampling account can therefore reproduce category-level
latency differences without classifying the faster category as intuitive.
This is a rival explanation requiring its own fit and measurements,
not an announcement that every two-process model is disproved.

Pennycook and colleagues' original reply distinguishes response-category
labeling from experimentally manipulating conflict. Both sequential
sampling and contemporary dual-process accounts can predict conflict
slowing. Their reply argues that response-specific effects of base-rate
extremity and metacognitive judgments may further constrain the accounts;
it does not claim every simple latency difference proves two processes.
The reply is a theoretical commentary on prior data, not a newly
recruited experiment \cite{PennycookEtAl2016RTCommentary}.

\subsection{Two responses, confidence, and individual differences}

The chapter reports that in some two-response tasks participants mostly
retain the first rapid answer when given an untimed second response.
This challenges a simple correction-only reading, but unchanged answers
also fit verification, failed reconsideration, or a stable choice of
strategy across trials. When both information sources are easy to
process, both proposals may already be available for the first answer.
Different problems produced more changes in Thompson and colleagues'
work; initial feeling-of-rightness (FOR) reports predicted subsequent
change and time. In an intervention account this is consistent with
monitoring influencing engagement, not a randomized demonstration that
FOR caused intervention \cite[p.~179]{Klauer2021DualProcess}.

Individual-differences correlations can also discriminate predictions
when specified in advance: capacity may matter more when resource-demanding
processing is needed than when an initial process already supplies the
task-correct response. They do not justify retroactively labeling every
capacity-correlated answer Type~2 or inferring that a lower-scoring
person has no such mechanism. Motivation, training, representation,
and task demands can contribute to the association.
The chapter reports that Stanovich and West found, across many tasks,
positive intelligence correlations with the designated normative response
and negative correlations with a modal response that departed from it.
It also reports confirmed instances of predictions about when such
correlations should and should not occur. No common coefficient, sample,
or universal benchmark is supplied for that collection of findings, and
the response that correlates with intelligence does not thereby acquire
normative authority \cite[p.~179]{Klauer2021DualProcess}.

\section{Dichotomous, gradual, decompositional, and relative views}

The source compares the dual-process debate with the history of
automaticity. A cluster definition is challenged when efficiency,
unintentionality, uncontrollability, and lack of consciousness fail to
co-occur. A decompositional response studies the features separately.
A gradual account retains a continuum: practice can replace a multi-step
algorithm by retrieval of an instance, or strengthen an algorithm.
Different features can change at different rates; lack of simultaneous
change alone does not refute that proposal. Opposite directions of
change with practice would be a more specific challenge
\cite[pp.~179--180]{Klauer2021DualProcess}.

An alternative elevates one feature to definitional status and treats
the others as correlates. Yet working-memory involvement itself comes
in degrees, and distinguishing very small from zero involvement is
difficult. Decoupling or representing reasons could supply a qualitative
criterion only after their meanings and measurements are specified.
A definition alone does not demonstrate two physically isolated systems.

The relative view introduces a comparison standard appropriate to a
task. In Stroop color naming, word reading is relatively more automatic
and interferes with naming an incongruent ink color; the converse
interference is not analogous when instructed to read. Both processes
may nevertheless count as relatively Type~1 when compared with a
demanding reasoning algorithm. Similarly, making belief assessment
harder and logical assessment easier can remove their usual relative
classification. This is an explicit context-dependent proposal, not
permission to change a classification merely to protect it from evidence.

Klauer's conclusion separates a framework's fruitfulness from a specific
model's success. The concepts have generated tasks, predictions, and
phenomena; particular accounts must still face graded single-process
and alternative-interaction explanations. A productive research program
is not proof of every architectural claim it has inspired.

\section{Source boundaries, comparisons, and usable tests}

Evans's Oxford chapter~8, \emph{Dual-Process Theories of Deductive
Reasoning: Facts and Fallacies}, is a separate 2012 source, pp.~115--133.
Its published abstract and contents identify criticisms of collapsing
all accounts together, positing exactly two systems, equating bias with
Type~1 and correctness with Type~2, treating context versus abstraction
as exhaustive, and inferring process type from speed. This comparison is
limited to that explicitly identified framing, not an assertion that
the two chapters are textually identical or that their detailed accounts
are interchangeable \cite{Evans2012DualProcessFactsFallacies}.
Klauer's four-account comparison, tripartite discussion, and relative-view
argument above retain their own 2021 source identity.

The linked handbook interfaces likewise stay separate:
\begin{itemize}
\item The heuristics chapter distinguishes environmental sampling and
      exposure from an internal procedure, and coherence from correspondence
      (\cref{sec:p19-rationality-2-4}). A biased output need not uniquely
      diagnose a Type~1 mechanism.
\item The mental-model chapter distinguishes an incomplete intuitive
      representation from explicit alternatives and counterexample search
      (\cref{sec:p19-rationality-2-3}). Representational format, computed
      function, and interaction architecture are three different commitments;
      a model-like representation does not identify all four accounts.
\item The monitoring chapter separates FOR, final confidence, object-level
      accuracy, and control (\cref{sec:p19-rationality-8-6}). Its primary
      CRT study is not a primary FOR experiment and cannot supply one.
\item Decision-theoretic metareasoning assigns costs and value to
      computation under stated performance laws (\cref{sec:p17-metareasoning}).
      A normative allocation policy is not the empirical reflective mind,
      and low deliberation is not irrational without goals and costs.
\end{itemize}

For a usable test, fix the account and rival before interpreting results;
separate logic, plausibility, resource load, instruction, and switching;
record the response and latency definitions and the participant/task
population; and measure initial representation or monitoring when those
are proposed mediators. A successful outcome is a prediction-sensitive
contrast under these conditions. A failed diagnosis is merely assigning
``intuitive'' to the observed fast or wrong category and ``analytic''
to the slow or correct one. The reviewed theories are empirical and
conceptual proposals, not established normative prescriptions to think
slowly in every situation.
